\documentclass[11pt,reqno]{amsart}

\usepackage[T1]{fontenc}
\usepackage{lmodern}
\usepackage{microtype}
\usepackage{mathtools}
\usepackage{amssymb}
\usepackage[colorlinks=true,linkcolor=blue,citecolor=blue,urlcolor=blue]{hyperref}

\numberwithin{equation}{section}

\newtheorem{theorem}{Theorem}[section]
\newtheorem{lemma}[theorem]{Lemma}
\newtheorem{corollary}[theorem]{Corollary}
\theoremstyle{remark}
\newtheorem{remark}[theorem]{Remark}

\newcommand{\Z}{\mathbb Z}
\newcommand{\Qp}{\mathbb Q_p}
\newcommand{\vp}{v_p}
\newcommand{\CT}{\operatorname {CT}}

\title[A multivariate Gauss supercongruence]
{A multivariate Gauss supercongruence for a Legendrian
Ap\'ery-like sequence}

\author{Anonymous}

\subjclass[2020]{Primary 11A07; Secondary 05A10, 11B65}
\keywords{Ap\'ery-like numbers, supercongruence, constant term,
multinomial coefficient, Jacobsthal congruence}

\begin{document}

\begin{abstract}
Let
\[
G_n=\sum_{k=0}^n4^k
\binom{2n-2k}{n-k}^{\!2}\binom{2k}{k}.
\]
Z.-H. Sun conjectured that
\[
G_{mp^r}\equiv G_{mp^{r-1}}\pmod {p^{2r}}
\]
for every odd prime \(p\) and all positive integers \(m,r\).
We prove a multivariate strengthening.  For
\[
\begin{split}
P(x,y)={}&(xy+x-y+1)(xy-x+y+1)\\
&{}\times(xy+x+y-1)(xy-x-y-1)
\end{split}
\]
and \(C(n;a,b)=[x^ay^b]P(x,y)^n\), we show that
\[
C(np^r;ap^r,bp^r)
\equiv C(np^{r-1};ap^{r-1},bp^{r-1})
\pmod {p^{2r}}.
\]
The proof separates the multinomial expansion into arrays divisible by
\(p\) and primitive arrays.  Two exponent-balance conditions force at least
two base-\(p\) carries at every relevant digit level in the primitive
sector.  The divisible sector is paired with the lower-level expansion
using the precise Jacobsthal--Kazandzidis congruence.  A constant-term
representation of \(G_n\) makes Sun's conjecture an immediate corollary.
It also gives the one-step congruence
\(G_N\equiv G_{N/p}\pmod {p^{2\vp(N)}}\) whenever \(p\) is an odd prime
dividing \(N\).
\end{abstract}

\maketitle

\section{Introduction}

Ap\'ery-like sequences sit at a meeting point of binomial identities,
differential equations, modular parametrizations, and arithmetic
congruences.  A recurring phenomenon is that their coefficients satisfy
congruences substantially stronger than those forced by an arbitrary
integral generating function.  Explaining the extra powers of a prime is
often the central difficulty.  We consider the second-order Legendrian
sequence in Zagier's classification \cite{Zagier2009},
\begin{equation}\label{eq:Gdef}
G_n=\sum_{k=0}^n4^k
\binom{2n-2k}{n-k}^{\!2}\binom{2k}{k}
\qquad(n\geq0).
\end{equation}
Its first terms are
\[
G_0=1,\qquad G_1=12,\qquad G_2=164,\qquad G_3=2352.
\]

In \cite[Conjecture~2.5]{Sun2020}, Z.-H. Sun conjectured that, for every
odd prime \(p\) and all positive integers \(m,r\),
\begin{equation}\label{eq:Sun}
G_{mp^r}\equiv G_{mp^{r-1}}\pmod {p^{2r}}.
\end{equation}
The modulus gains two powers of \(p\) at every dilation.  This is the feature
that needs explanation: a first-order Gauss congruence would account for only
one new power.  Our proof identifies a finite combinatorial mechanism for the
second power.  Gorodetsky's two-variable constant-term representation
\cite[Sec.~5]{GorodetskyV2}, after a quadratic substitution, factors into
four signed linear forms.  In the associated multinomial expansion, the two
independent exponent balances force two base-\(p\) carries at each relevant
digit level.  The primitive terms therefore already contain the full modulus;
the remaining terms pair with the preceding level by the
Jacobsthal--Kazandzidis congruence.

This viewpoint places the result near, but not inside, two established
frameworks.  The general Dwork congruences of Mellit and Vlasenko
\cite{MellitVlasenko} naturally give first-order congruences when the Newton
polytope has a unique interior lattice point; they do not formally supply the
extra power needed here.  Straub \cite{Straub2014} proved second-order, and
in a subfamily third-order, multivariate supercongruences for a
partition-indexed family of rational functions.  The rational function
\(1/(1-z\Lambda_G(x,y))\) below is not a displayed member of that family,
and our balanced four-row coefficient model has a different combinatorial
source.  Thus neither result gives \eqref{eq:Sun} by specialization.

The natural statement revealed by the constant-term representation is
stronger than Sun's diagonal congruence.  Put

\begin{equation}\label{eq:Pdef}
\begin{split}
P(x,y)={}&(xy+x-y+1)(xy-x+y+1)\\
&{}\times(xy+x+y-1)(xy-x-y-1)
\end{split}
\end{equation}
and define
\[
C(n;a,b)=[x^ay^b]P(x,y)^n
\qquad(n,a,b\geq0).
\]
For later use, we set \(C(n;a,b)=0\) if \(a<0\) or \(b<0\).
Our main result is the following.

\begin{theorem}\label{thm:multi}
Let \(p\) be an odd prime, \(r\geq1\), and \(n,a,b\geq0\).  Then
\begin{equation}\label{eq:multi}
C(np^r;ap^r,bp^r)
\equiv C(np^{r-1};ap^{r-1},bp^{r-1})
\pmod {p^{2r}}.
\end{equation}
\end{theorem}

Theorem~\ref{thm:multi} proves more than a diagonal congruence: its Laurent
form controls every coefficient on a simultaneously dilated lattice.  It is
an order-two coefficient congruence for the Legendrian polynomial occurring
in \cite[(5.2)]{GorodetskyV2}, analogous in spirit to the questions raised in
\cite[(6.2)]{GorodetskyV2} but not covered by the list in that question.

\begin{corollary}\label{cor:Sun}
Z.-H. Sun's congruence \eqref{eq:Sun} holds for every odd prime \(p\) and
all positive integers \(m,r\).
\end{corollary}

The proof has two complementary halves.  For an array in the multinomial
expansion that is not divisible by \(p\), Legendre's formula converts its
valuation into a sum of carry counts.  The row sums and the two exponent
balances rule out a single carry, giving at least two at every digit level.
For an array divisible by \(p\), division by \(p\) produces exactly one term
at the preceding level; the Jacobsthal--Kazandzidis quotient supplies the
remaining precision after the common \(p\)-power has been removed.  The two
valuation contributions add to \(2r\) independently of the depth.  This is
why no separate induction on exceptional low layers is needed.

Section~2 derives the constant-term model directly from the defining
convolution.  Section~3 proves the two-carry lemma, Section~4 records the
scaled multinomial estimate, and Section~5 combines them into the
multivariate theorem.  Section~6 returns to the Laurent polynomial and
deduces Sun's conjecture.  The final section records exact computations used
only as checks; the proof itself is independent of computation.

\section{The constant-term representation}

We first derive the representation directly from \eqref{eq:Gdef}.  Set
\[
B(t)=\sum_{n\geq0}\binom{2n}{n}^{\!2}t^n
={}_2F_1\left(\frac12,\frac12;1;16t\right).
\]
The convolution in \eqref{eq:Gdef} gives
\begin{equation}\label{eq:GFfirst}
\sum_{n\geq0}G_nt^n=\frac{B(t)}{\sqrt{1-16t}}.
\end{equation}
Pfaff's transformation yields
\[
B(t)=\frac1{\sqrt{1-16t}}\,
B\left(\frac{-t}{1-16t}\right).
\]
Substitution in \eqref{eq:GFfirst} and extraction of the coefficient of
\(t^n\) give
\begin{equation}\label{eq:alternating}
G_n=\sum_{k=0}^n(-1)^k16^{n-k}
\binom nk\binom{2k}{k}^{\!2}.
\end{equation}
Since
\[
\CT_x(x+2+x^{-1})^k=\binom{2k}{k},
\]
equation \eqref{eq:alternating} is equivalent to
\begin{equation}\label{eq:CTsimple}
G_n=\CT_{x,y}
\left(16-(x+2+x^{-1})(y+2+y^{-1})\right)^n.
\end{equation}

Changing \(x,y\) to \(-x,-y\) in \eqref{eq:CTsimple} gives the Laurent
polynomial recorded in \cite[(5.2)]{GorodetskyV2}:
\begin{equation}\label{eq:LambdaG}
\Lambda_G(x,y)=(-xy)^{-1}
\left(x^2y^2-2x^2y-2xy^2+x^2-12xy+y^2-2x-2y+1\right).
\end{equation}
Thus
\[
G_n=\CT\Lambda_G(x,y)^n.
\]

Let
\[
A=(x-x^{-1})(y-y^{-1}).
\]
A direct calculation gives
\[
\Lambda_G(x^2,y^2)=16-A^2=-\frac{P(x,y)}{x^2y^2},
\]
because
\[
\begin{split}
xy(4+A)&=(xy+x-y+1)(xy-x+y+1),\\
xy(4-A)&=-(xy+x+y-1)(xy-x-y-1).
\end{split}
\]
Substitution by \(x^2,y^2\) does not change a constant term.  Consequently
\begin{equation}\label{eq:Gcoefficient}
G_n=(-1)^n[x^{2n}y^{2n}]P(x,y)^n
=(-1)^nC(n;2n,2n).
\end{equation}

\section{Balanced arrays and two forced carries}

Fix nonnegative integers \(N,A,B\).  Expand each of the four factors in
\eqref{eq:Pdef} to the \(N\)th power.  For the \(i\)th factor let
\[
q_i=(a_i,b_i,c_i,d_i),\qquad
a_i+b_i+c_i+d_i=N,
\]
where the four entries record the multiplicities of \(xy,x,y,1\),
respectively.  A four-row array
\[
q=(q_1,q_2,q_3,q_4)
\]
contributes to \(C(N;A,B)\) precisely when
\begin{equation}\label{eq:balances}
\sum_{i=1}^4(a_i+b_i)=A,\qquad
\sum_{i=1}^4(a_i+c_i)=B.
\end{equation}
We call such an array admissible.  The absolute value of its contribution is
\begin{equation}\label{eq:PN}
P_N(q)=\prod_{i=1}^4
\binom{N}{a_i,b_i,c_i,d_i}.
\end{equation}
The sign is a parity character in the entries of \(q\).  Explicitly, the
four factors in \eqref{eq:Pdef} contribute the sign exponents
\[
c_1,\qquad b_2,\qquad d_3,\qquad b_4+c_4+d_4.
\]
It follows that multiplication of all entries by an odd integer preserves
the sign.

The next lemma is the source of the exponent \(2R\).  At a fixed base-\(p\)
level, the residue mass of a primitive array cannot be concentrated in a
single carry: the two balances would force that mass into one column, while
the row sums would force it into one row, contradicting the choice of least
residues.  Thus every level contributes at least two carries.  We give the
full argument because this small rigidity statement is the only place where
the special four-factor geometry of \(P\) enters the valuation estimate.

\begin{lemma}\label{lem:primitive}
Let \(p\) be a prime, let \(R\geq0\), and suppose that \(p^R\) divides
each of \(N,A,B\).
If \(q\) is an admissible array not all of whose entries are divisible by
\(p\), then
\[
\vp(P_N(q))\geq2R.
\]
\end{lemma}

\begin{proof}
The assertion is immediate when \(R=0\), so assume \(R\geq1\).
Fix \(1\leq j\leq R\) and put \(M=p^j\).  Replace each of the sixteen
entries of \(q\) by its least residue in \([0,M-1]\).  Let the four column
totals of these residues be \(X,Y,Z,W\), in the order corresponding to
\(a,b,c,d\), and write
\[
X+Y+Z+W=MT_j.
\]
Each row sum is congruent to \(N\), hence to zero, modulo \(M\), so
\(T_j\) is a nonnegative integer.  It is positive because at least one
entry of \(q\) is not divisible by \(p\).

We show that \(T_j\neq1\).  Suppose otherwise.  Then the total residue mass
is \(M\).  The two balances \eqref{eq:balances} give
\[
X+Y\equiv X+Z\equiv0\pmod M.
\]
Both \(X+Y\) and \(X+Z\) lie between \(0\) and \(M\), and hence each is
either \(0\) or \(M\).  Considering the four possible pairs of values shows
that all residue mass must lie in one column:
\[
(X,Y,Z,W)\in
\{(M,0,0,0),(0,M,0,0),(0,0,M,0),(0,0,0,M)\}.
\]
On the other hand, every row residue sum is a multiple of \(M\), and the
total of the four row sums is \(M\).  Thus exactly one row has residue sum
\(M\).  Since the mass lies in a single column, one least residue would
equal \(M\), a contradiction.  Therefore \(T_j\geq2\).

More explicitly, Legendre's formula gives the level-\(j\) contribution to
the valuation of the product in \eqref{eq:PN} as
\[
 \sum_{i=1}^4\left(\left\lfloor\frac{N}{M}\right\rfloor
 -\sum_{\alpha\in\{a,b,c,d\}}
  \left\lfloor\frac{q_{i,\alpha}}{M}\right\rfloor\right)
 =\frac{X+Y+Z+W}{M}=T_j.
\]
All contributions for \(j>R\) are nonnegative.  Hence
\[
\vp(P_N(q))\geq\sum_{j=1}^RT_j\geq2R.
\]
\end{proof}

\begin{remark}
The lower bound is sharp.  When \(N=p^R\) and \(A=B=2N\), the rows
\[
(0,0,0,N),\quad(0,0,1,N-1),\quad
(N-1,1,0,0),\quad(N,0,0,0)
\]
form an admissible array with total valuation \(2R\).
\end{remark}

\section{Scaled multinomial coefficients}

We use the following consequence of the Jacobsthal--Kazandzidis
congruence.  The quotient congruences in this section are interpreted in
\(\Qp\): thus \(X\equiv Y\pmod {p^e}\) means
\(\vp(X-Y)\geq e\).

\begin{lemma}\label{lem:Jacobsthal}
Let \(p\) be an odd prime, \(s\geq0\), and let
\(\boldsymbol{v}=(v_1,\ldots,v_t)\) be a weak composition (zero parts are
allowed).  Then
\begin{equation}\label{eq:multiJacob}
\frac{\binom{p^{s+1}|\boldsymbol v|}
 {p^{s+1}v_1,\ldots,p^{s+1}v_t}}
 {\binom{p^s|\boldsymbol v|}
 {p^sv_1,\ldots,p^sv_t}}
\equiv1\pmod {p^{2s+2}}.
\end{equation}
The numerator and denominator in \eqref{eq:multiJacob} have the same
\(p\)-adic valuation.
\end{lemma}

\begin{proof}
Decompose each multinomial coefficient into \(t-1\) binomial coefficients.
Trivial factors produced by zero parts are omitted.
For a nontrivial resulting factor, apply the classical congruence
\begin{equation}\label{eq:Jacob}
\frac{\binom{pn}{pm}}{\binom nm}
\equiv1\pmod {p^{\vp(nm(n-m))+3-\delta_{p,3}}},
\end{equation}
where \(\delta_{p,3}=1\) for \(p=3\) and zero otherwise; see
\cite{Jacobsthal1952}, \cite[p.~255]{Granville1997}, or
\cite[Lemma~4.1]{Gorodetsky2023}.
In the present application, \(n,m,n-m\) contain the common factor \(p^s\).
The exponent in \eqref{eq:Jacob} is therefore at least \(3s+2\) for \(p=3\)
and at least \(3s+3\) for \(p\geq5\).  Both bounds are at least \(2s+2\),
so multiplication proves \eqref{eq:multiJacob}.

The valuation assertion follows from Legendre's digit-sum formula:
multiplication of all arguments by \(p\) merely shifts their base-\(p\)
digits.
\end{proof}

\section{Proof of the multivariate congruence}

We prove a slightly stronger one-step statement.

\begin{theorem}\label{thm:strong}
Let \(p\) be an odd prime, and let \(N,A,B\geq0\) be not all zero.
Suppose that \(p\) divides each of \(N,A,B\), and put
\[
R=\min\{\vp(N),\vp(A),\vp(B)\},
\]
with the convention \(\vp(0)=+\infty\).  Then
\begin{equation}\label{eq:strong}
C(N;A,B)\equiv C(N/p;A/p,B/p)\pmod {p^{2R}}.
\end{equation}
\end{theorem}

\begin{proof}
If \(N=0\), then \(A+B>0\), and both coefficients in
\eqref{eq:strong} are zero.  We may therefore assume \(N>0\); in
particular, \(R\geq1\).
Split the admissible arrays contributing to \(C(N;A,B)\) into two classes.
By Lemma~\ref{lem:primitive}, every array not wholly divisible by \(p\)
contributes zero modulo \(p^{2R}\).

It remains to compare the wholly divisible arrays term by term with the
lower-level expansion.  The following bookkeeping is arranged so that the
valuation already present in the lower term and the precision of its scaling
quotient add to exactly \(2R\).

Every array in the other class is uniquely \(q=pu\), where \(u\) is
admissible for \(C(N/p;A/p,B/p)\).  This gives a bijection with all arrays
in the lower expansion, and the signs agree because \(p\) is odd.

Let
\[
s=\min_{i,\alpha}\vp(u_{i,\alpha}),
\]
where zero entries have infinite valuation, and write \(u=p^sv\).
At least one of \(N/p,A/p,B/p\) has exact valuation \(R-1\), so
\(0\leq s\leq R-1\).  The array \(v\) is primitive, and its three
parameters are all divisible by \(p^{R-1-s}\).  Common scaling does not
change a multinomial valuation.  Lemma~\ref{lem:primitive} therefore gives
\begin{equation}\label{eq:lowerVal}
\vp(P_{N/p}(u))\geq2(R-1-s).
\end{equation}

Apply Lemma~\ref{lem:Jacobsthal} to each of the four rows and multiply the
four \(p\)-adic unit quotients.  We obtain
\begin{equation}\label{eq:ratio}
\frac{P_N(pu)}{P_{N/p}(u)}
\equiv1\pmod {p^{2s+2}}.
\end{equation}
Equations \eqref{eq:lowerVal} and \eqref{eq:ratio} imply
\[
\vp\bigl(P_N(pu)-P_{N/p}(u)\bigr)
\geq2(R-1-s)+2s+2=2R.
\]
Thus every paired divisible term agrees modulo \(p^{2R}\).  Adding these
terms and the primitive terms proves \eqref{eq:strong}.
\end{proof}

\begin{proof}[Proof of Theorem~\ref{thm:multi}]
Apply Theorem~\ref{thm:strong} with
\[
(N,A,B)=(np^r,ap^r,bp^r).
\]
If \(n=a=b=0\), both coefficients in \eqref{eq:multi} are one; otherwise
the minimum valuation is at least \(r\).
\end{proof}

\section{Laurent form and Sun's conjecture}

Set
\[
\Lambda(x,y)=-\frac{P(x,y)}{x^2y^2}.
\]
For \(n\geq0\) and \(h,k\in\Z\), define
\[
D(n;h,k)=\CT\bigl(\Lambda(x,y)^n x^hy^k\bigr).
\]
Theorem~\ref{thm:multi} is equivalent to
\begin{equation}\label{eq:LaurentMulti}
D(np^r;hp^r,kp^r)
\equiv D(np^{r-1};hp^{r-1},kp^{r-1})
\pmod {p^{2r}}.
\end{equation}
Indeed,
\[
D(n;h,k)=(-1)^nC(n;2n-h,2n-k),
\]
where the convention following the definition of \(C\) is used.  The two
signs agree under dilation by an odd prime.  If both ordinary-polynomial
indices are nonnegative, Theorem~\ref{thm:multi} applies directly, whether
or not the coefficient happens to vanish.  If either index is negative,
the corresponding index remains negative after dilation, so both sides of
\eqref{eq:LaurentMulti} are zero.

\begin{proof}[Proof of Corollary~\ref{cor:Sun}]
By \eqref{eq:Gcoefficient},
\[
G_n=(-1)^nC(n;2n,2n).
\]
Apply Theorem~\ref{thm:multi} with
\[
(n,a,b)=(m,2m,2m).
\]
The signs at levels \(r\) and \(r-1\) agree because \(p\) is odd, and
\eqref{eq:Sun} follows.
\end{proof}

\begin{remark}
Theorem~\ref{thm:strong} also yields the stronger one-step form
\[
G_N\equiv G_{N/p}\pmod {p^{2\vp(N)}}
\]
for every odd prime \(p\mid N\).
\end{remark}

\section{Computational checks}

The proof above is independent of computation.  Exact integer programs were
nevertheless used to check the transformations and boundary cases.
The original convolution \eqref{eq:Gdef}, the alternating formula
\eqref{eq:alternating}, and the signed four-factor expansion were computed
through independent code paths.  A min-plus enumeration found the sharp
minimum \(2R\) in five primitive-sector regimes.  Finally, all scaled
coefficient pairs were checked in the eight regimes
\[
\begin{split}
(p,r,n)\in\{&(3,1,1),(3,1,2),(3,2,1),(3,2,2),\\
&(5,1,1),(5,1,2),(5,2,1),(7,1,1)\}.
\end{split}
\]
All 368 coefficient congruences passed.
The four exact verification programs used for these checks are supplied in
the public repository accompanying this manuscript; they are not part of the
logical proof.

\newpage
\section*{Statements and declarations}

\noindent\textbf{Funding.}
No funding was received to assist with the preparation of this manuscript.

\smallskip
\noindent\textbf{Competing interests.}
The author has no relevant financial or non-financial interests to disclose.

\smallskip
\noindent\textbf{Data availability.}
No datasets were generated or analysed during the current study.

\smallskip
\noindent\textbf{Code availability.}
The exact verification programs are provided at
the \href{https://github.com/huiminZheng-collab/sun-conjecture-2-5-proof-candidate}
{public GitHub repository}.

\smallskip
\noindent\textbf{Use of generative AI.}
The proof presented in this article was found by OpenAI Codex.

\footnotesize
\begin{thebibliography}{99}

\bibitem{Gorodetsky2023}
O. Gorodetsky,
\emph{New representations for all sporadic Ap\'ery-like sequences, with
applications to congruences},
Exp. Math. \textbf{32} (2023), no.~4, 641--656.
\url{https://doi.org/10.1080/10586458.2021.1982080}.

\bibitem{GorodetskyV2}
O. Gorodetsky,
\emph{New representations for all sporadic Ap\'ery-like sequences, with
applications to congruences},
arXiv:2102.11839v2 (2025).
\url{https://arxiv.org/abs/2102.11839v2}.

\bibitem{Granville1997}
A. Granville,
\emph{Arithmetic properties of binomial coefficients. I. Binomial
coefficients modulo prime powers},
in: Organic Mathematics (Burnaby, BC, 1995), CMS Conf. Proc. \textbf{20},
Amer. Math. Soc., Providence, RI, 1997, 253--276.

\bibitem{Jacobsthal1952}
V. Brun, J. O. Stubban, J. E. Fjeldstad, R. Tambs Lyche, K. E. Aubert,
W. Ljunggren and E. Jacobsthal,
\emph{On the divisibility of the difference between two binomial
coefficients},
in: Den 11te Skandinaviske Matematikerkongress, Trondheim, 1949,
Johan Grundt Tanums Forlag, Oslo, 1952, 42--54.

\bibitem{MellitVlasenko}
A. Mellit and M. Vlasenko,
\emph{Dwork's congruences for the constant terms of powers of a Laurent
polynomial},
Int. J. Number Theory \textbf{12} (2016), no.~2, 313--321.
\url{https://doi.org/10.1142/S1793042116500184}.

\bibitem{Straub2014}
A. Straub,
\emph{Multivariate Ap\'ery numbers and supercongruences of rational
functions},
Algebra Number Theory \textbf{8} (2014), no.~8, 1985--2008.
\url{https://doi.org/10.2140/ant.2014.8.1985}.

\bibitem{Sun2020}
Z.-H. Sun,
\emph{New congruences involving Ap\'ery-like numbers},
arXiv:2004.07172v2 (2020).
\url{https://arxiv.org/abs/2004.07172v2}.

\bibitem{Zagier2009}
D. Zagier,
\emph{Integral solutions of Ap\'ery-like recurrence equations},
in: Groups and Symmetries, CRM Proc. Lecture Notes \textbf{47},
Amer. Math. Soc., Providence, RI, 2009, 349--366.

\end{thebibliography}
\normalsize

\end{document}
